\documentclass[11pt]{article}
\usepackage[T1]{fontenc}
\usepackage{lmodern}
\usepackage[margin=1in]{geometry}
\usepackage{amsmath,amssymb,amsthm}
\usepackage{microtype}
\usepackage[hidelinks]{hyperref}

\newtheorem*{sourcequestion}{Source question}
\newtheorem*{answertheorem}{Later theorem}

\title{Two Random-Forcing Questions from arXiv:1909.09387\\Answered by arXiv:2110.04568}
\author{Literature identification note}
\date{August 17, 2026}

\begin{document}
\maketitle

\begin{abstract}
Two questions in Sobota and Zdomskyy's \emph{Convergence of measures in
forcing extensions} are explicitly answered by their later paper
\emph{Convergence of measures after adding a real}.  Random forcing destroys
both the Nikodym and Grothendieck properties of every infinite ground-model
Boolean algebra, and the random model satisfies
$\mathfrak d<\mathfrak{nik}=\mathfrak{gr}$.  These results settle source
Questions~6.13 and~6.4, respectively.
\end{abstract}

\section*{Identification status}

\begin{itemize}
\item Run \texttt{fa\_banach\_001}, lane 13.
\item Source: arXiv:1909.09387 \cite{source}.
\item Answer: arXiv:2110.04568 \cite{answer}.
\item Classification: explicit later-literature answer; no new theorem claimed.
\end{itemize}

\section{Random forcing and the Vitali--Hahn--Saks property}

\begin{sourcequestion}[Question 6.13 of \cite{source}, PDF page 20]
Let $\mathbb P$ be random forcing and let $\mathcal A$ be a ground-model
Boolean algebra which is $\sigma$-complete, or merely has property~(I).
Must $\mathcal A$ have the Vitali--Hahn--Saks property in the generic
extension?
\end{sourcequestion}

\begin{answertheorem}[Theorem 4.6 of \cite{answer}, PDF page 10]
If a forcing adds a random real, then every infinite ground-model Boolean
algebra has neither the Nikodym property nor the Grothendieck property in the
generic extension.
\end{answertheorem}

The answer to Question~6.13 is therefore \emph{no}, in the stronger form that
neither $\sigma$-completeness nor property~(I) is relevant.  Indeed, the
Vitali--Hahn--Saks property implies both the Nikodym and Grothendieck
properties, whereas the later theorem destroys both for every infinite
ground-model algebra.

\section{The cardinal invariants above \texorpdfstring{$\mathfrak d$}{d}}

\begin{sourcequestion}[Question 6.4 of \cite{source}, PDF page 18]
For $\mathfrak x\in
\{\mathfrak{nik},\mathfrak{gr},\mathfrak{vhs}\}$, is it consistent that
$\mathfrak x>\mathfrak d$?  In particular, is there an
$\omega^\omega$-bounding forcing after which
$\mathcal P(\omega)^V$ loses the Vitali--Hahn--Saks property?
\end{sourcequestion}

Corollary~5.3 of \cite{answer} (PDF page~20) states
\[
 \mathfrak d<\mathfrak{nik}=\mathfrak{gr}
 \qquad\text{in the random model}.
\]
The source records
$\mathfrak{nik},\mathfrak{gr}\le\mathfrak{vhs}$, so this same model has
$\mathfrak d<\mathfrak{vhs}$ as well.  Thus Question~6.4 has a positive
answer for every listed $\mathfrak x$.  Moreover, measure-algebra random
forcing is $\omega^\omega$-bounding, and Theorem~4.6 applied to the infinite
algebra $\mathcal P(\omega)^V$ supplies exactly the forcing requested in the
particular clause.

\section*{Scope}

The later paper is by the same authors and explicitly describes the random
case as addressing the preservation issue left open in the earlier work.
This identification settles Questions~6.4 and~6.13 only; the other questions
in Section~6 of arXiv:1909.09387 are not classified here.

\begin{thebibliography}{9}
\bibitem{source}
D.~Sobota and L.~Zdomskyy,
\emph{Convergence of measures in forcing extensions},
arXiv:1909.09387 (2019).

\bibitem{answer}
D.~Sobota and L.~Zdomskyy,
\emph{Convergence of measures after adding a real},
arXiv:2110.04568 (2021; published 2024).
\end{thebibliography}

\end{document}
